\documentclass[11pt,a4paper]{article}

\usepackage[T1]{fontenc}
\usepackage{lmodern}
\usepackage[a4paper,margin=21mm,headheight=15pt]{geometry}
\usepackage{microtype}
\usepackage{amsmath,amssymb}
\usepackage{booktabs}
\usepackage{enumitem}
\usepackage[numbers,sort&compress]{natbib}
\usepackage{graphicx,float,caption}
\usepackage{xcolor}
\usepackage{fancyhdr}
\usepackage{hyperref}

\definecolor{reportink}{HTML}{163044}
\definecolor{reportaccent}{HTML}{12615E}
\hypersetup{
  colorlinks=true, linkcolor=reportaccent, citecolor=reportaccent, urlcolor=reportaccent,
  pdftitle={Neuronal firing and excitability with the FitzHugh--Nagumo model},
  pdfauthor={Mathematical Modelling Group Project (MAC F314)}
}
\captionsetup{font=small,labelfont={bf,color=reportink},hypcap=false}
\setlength{\parindent}{0pt}
\setlength{\parskip}{0.45em}
\setlength{\emergencystretch}{1em}
\setlength{\tabcolsep}{7pt}

\pagestyle{fancy}
\fancyhf{}
\fancyhead[L]{\small\color{reportink} MAC F314 \enspace|\enspace Mathematical Modelling}
\fancyhead[R]{\small\color{reportink} FitzHugh--Nagumo model}
\fancyfoot[C]{\small\color{reportink}\thepage}
\renewcommand{\headrulewidth}{0.4pt}

% ---------------------------------------------------------------------------
%  GROUP MEMBERS.  Edit only here; the title block reads from these macros.
% ---------------------------------------------------------------------------
\newcommand{\memberA}{Bhagyesh Kamal}
\newcommand{\memberB}{Jaivin Bengani}
\newcommand{\memberC}{Purunjay Karwa}
\newcommand{\memberD}{Arnav Kulshreshtha}
\newcommand{\memberE}{Soham Agrawal}
\newcommand{\memberF}{Samar Nayak}

\begin{document}

{\color{reportaccent}\small\bfseries MATHEMATICAL MODELLING GROUP PROJECT \enspace MAC F314\par}
\vspace{0.45em}
{\color{reportink}\LARGE\bfseries Neuronal firing and excitability\\[0.15em]
with the FitzHugh--Nagumo model\par}
\vspace{0.35em}
{\large Equilibria, bifurcations, numerical simulation, and sensitivity\par}
\vspace{0.5em}
{\small
\begin{tabular}{@{}lll@{}}
\memberA & \memberC & \memberE\\
\memberB & \memberD & \memberF\\
\end{tabular}\par}

\vspace{0.7em}
\noindent\textcolor{reportaccent}{\rule{\linewidth}{1pt}}

\subsection*{Abstract}

We use the FitzHugh--Nagumo model to reproduce rest, a single spike, and repeated firing. A positive first Lyapunov coefficient classifies two current crossings, near \(I=0.331\) and \(I=1.419\), as subcritical Hopf bifurcations. At \(I=0.50\), raising the recovery rate from \(0.06\) to \(0.10\) shortens the period from \(49.157\) to \(33.525\) model time units. In the controlled inverse study, least squares gave lower held-out voltage error than the PINN in all ten noisy comparisons.

\section{Introduction and motivation}

A brief electrical disturbance can produce a voltage spike followed by recovery. A suitable sustained current can produce repeated firing. The FitzHugh--Nagumo model represents both responses with fast voltage and slower recovery variables. Its two-dimensional phase plane lets us relate simulated trajectories to equilibrium and stability calculations \cite{fitzhugh1961,nagumo1962}. Excitable systems of this form are treated as a standard topic in mathematical biology by Edelstein-Keshet \cite{edelsteinkeshet2005}, and the modelling strategy of reducing a detailed mechanism to the smallest system that still reproduces the behaviour of interest follows Dym \cite{dym2004}.

\subsection*{Objectives}

\begin{enumerate}
  \item Derive the nullclines and equilibrium equation, and determine local stability from the Jacobian.
  \item Use controlled simulations to compare return to rest, a single spike, and repetitive firing.
  \item Locate and classify the bifurcations as current varies, and measure how the recovery parameter affects the interval between spikes.
\end{enumerate}

\section{Assumptions and mathematical formulation}

We model one spatially uniform excitable cell. The state and time are scaled, and the simulations are deterministic. Applied current is constant within each run. A brief nudge is represented as an initial change in voltage while the starting recovery value is fixed. The baseline parameters are \(a=0.7\), \(b=0.8\), and \(\varepsilon=0.08\).

\begin{align}
  \dot v &= v-\frac{v^3}{3}-w+I, \label{eq:voltage}\\
  \dot w &= \varepsilon(v+a-bw). \label{eq:recovery}
\end{align}

Here \(v\) is scaled voltage, \(w\) is recovery, \(I\) is applied current, and a dot denotes a time derivative. Higher recovery lowers the voltage slope through \(-w\), while higher voltage raises the recovery slope. The small factor \(\varepsilon\) slows recovery on the time scale of a spike.

\section{Mathematical analysis}

\subsection{Nullclines and equilibria}

A nullcline is a curve on which one rate of change is zero. Setting \(\dot v=0\) gives the cubic voltage nullcline. Setting \(\dot w=0\) gives the straight recovery nullcline.
\begin{equation}
  w=v-\frac{v^3}{3}+I,
  \qquad
  w=\frac{v+a}{b}.
  \label{eq:nullclines}
\end{equation}
An equilibrium is an intersection of these curves, where both rates vanish. Equating the two expressions for \(w\) gives
\begin{equation}
  \frac{b}{3}v_*^3+(1-b)v_*+a-bI=0,
  \qquad
  w_*=\frac{v_*+a}{b}.
  \label{eq:equilibrium}
\end{equation}
The derivative of the cubic with respect to \(v_*\) is \(bv_*^2+1-b>0\) for \(b=0.8\). The cubic is strictly increasing and takes arbitrarily large positive and negative values, so there is exactly one equilibrium at each current. Increasing \(I\) shifts the voltage nullcline upward and moves the equilibrium to a higher voltage.

\subsection{Jacobian and local stability}

The Jacobian records how the two rates respond to small changes in voltage and recovery. At the equilibrium it is
\begin{equation}
  J(v_*)=
  \begin{pmatrix}
    1-v_*^2 & -1\\
    \varepsilon & -\varepsilon b
  \end{pmatrix}.
\end{equation}
Its trace and determinant are
\begin{equation}
  T=1-v_*^2-\varepsilon b,
  \qquad
  D=\varepsilon\bigl[1-b(1-v_*^2)\bigr].
  \label{eq:trace-det}
\end{equation}
The eigenvalues solve \(\lambda^2-T\lambda+D=0\). Negative real parts imply local asymptotic stability. A positive real part implies instability. With our parameters, \(D=\varepsilon(1-b+bv_*^2)>0\), so away from \(T=0\) the sign of the trace determines whether the equilibrium is stable or unstable. The Jacobian predicts the initial response to a small nudge. A larger nudge can explore regions of the phase plane where the linear approximation is inaccurate.

\subsection{A numerically checked region around the resting state}
\label{sec:lyapunov}

The Jacobian establishes that \emph{sufficiently small} disturbances decay, without saying how small. A quadratic Lyapunov function lets us attach a number to it. At the zero-current equilibrium $A=J(v_*)$ is Hurwitz, so the Lyapunov matrix equation $A^{\mathsf{T}}P+PA=-Q$ with $Q=\mathbb{I}$ has a unique positive-definite solution,
\begin{equation}
  P=\begin{pmatrix}1.091492 & -0.266174\\ -0.266174 & 11.971474\end{pmatrix},
  \qquad \lambda(P)=\{1.08498,\;11.97798\}.
  \label{eq:lyapP}
\end{equation}
Write $z=(v-v_*,\,w-w_*)^{\mathsf{T}}$ and $V(z)=z^{\mathsf{T}}Pz$. Differentiating along the \emph{full nonlinear} system $\dot z=Az+g(z)$, with $g(z)=(-v_*z_1^2-z_1^3/3,\,0)^{\mathsf{T}}$ collecting the quadratic and cubic terms,
\begin{equation}
  \dot V=-\lVert z\rVert^2+2\,(Pz)_1\Bigl(-v_*z_1^2-\tfrac{z_1^3}{3}\Bigr).
  \label{eq:Vdot}
\end{equation}
If $\dot V$ were negative on the whole of a sublevel set $\{V\le c\}$ away from the origin, that set would be forward invariant and contained in the basin of attraction. We do not establish that analytically here. Instead we search outward along rays for the level at which $\dot V$ first reaches zero, step back to a level below it, and sample. On
\begin{equation}
  \Omega=\{z:V(z)\le 0.2059\},
  \label{eq:lyapset}
\end{equation}
evaluating $\dot V$ at $4\times10^5$ points of the boundary and on $400$ interior level sets, the largest value found is $-4.05\times10^{-6}$. We therefore report $\Omega$ as a \emph{numerically checked region}: at every point sampled, $\dot V$ is negative. Along the voltage axis, the numerically checked region reaches a displacement of $0.43433$, which is $78.2\%$ of the displacement $0.5555$ at which the chosen spike criterion is met (Section~\ref{sec:threshold}). The $+0.05$ displacement used in Section~\ref{sec:runs} has $V=0.002729$, well inside $\Omega$.

The limits of this are worth stating precisely. A grid establishes the sign of $\dot V$ at the points sampled; it does not establish it between them, so we do not claim that $\dot V<0$ everywhere in $\Omega$, nor that every trajectory starting in $\Omega$ returns to rest. The margin $-4.05\times10^{-6}$ is small, and the level is a numerical estimate rather than a proof that no larger one works. Converting this into a guarantee needs a bound on $\dot V$ between sample points --- a Lipschitz estimate on $\dot V$ over $\Omega$ combined with the sample spacing would do it --- and that is the natural next step.


\clearpage

\subsection{Locating and classifying the Hopf bifurcations}

When \(T=0\) and \(D>0\), the eigenvalues are purely imaginary. Solving \(T=0\) and substituting into \eqref{eq:nullclines} gives
\begin{equation}
  v_H=\pm\sqrt{1-\varepsilon b},
  \qquad
  I_H=\frac{v_H+a}{b}-v_H+\frac{v_H^3}{3}.
  \label{eq:hopf}
\end{equation}
At both candidates \(D=0.075904\). Differentiating \eqref{eq:equilibrium} with respect to \(I\) and using \eqref{eq:trace-det} gives
\begin{equation}
  \frac{dT}{dI}=-\frac{2bv_*}{bv_*^2+1-b}.
\end{equation}
The nonzero slopes show that the eigenvalue pair crosses the imaginary axis. To classify the bifurcations we also need the effect of the nonlinear terms on the amplitude of small oscillations \cite{kuznetsov2006}.

\subsubsection*{First Lyapunov coefficient}

At either crossing, set \(x=v-v_H\), \(y=w-w_H\), \(p=\varepsilon b\), and \(\omega=\sqrt{\varepsilon-p^2}\). Subtract the equilibrium equations and expand the cubic to obtain
\[
  \dot x=px-y+F(x),\qquad
  \dot y=\varepsilon x-py,\qquad
  F(x)=-v_Hx^2-\frac{x^3}{3}.
\]
Use \(X=x\), \(Y=(y-px)/\omega\), and \(\tau=\omega t\) to put the linear part into unit-frequency rotation form. With primes denoting \(\tau\) derivatives,
\[
  X'=-Y+f(X),\qquad Y'=X+g(X),\qquad
  f=F/\omega,\quad g=-pF/\omega^2.
\]
The planar coefficient formula \cite{guckenheimer1983} simplifies because \(f\) and \(g\) depend only on \(X\). At the origin, \(f_{XXX}=-2/\omega\), \(f_{XX}=-2v_H/\omega\), and \(g_{XX}=2pv_H/\omega^2\). Substitution gives
\begin{equation}
  \ell_1=\frac{f_{XXX}-f_{XX}g_{XX}}{16}
        =-\frac{1}{8\omega}+\frac{pv_H^2}{4\omega^3}
        =\frac{\varepsilon(2b-1-\varepsilon b^2)}{8\omega^3}.
  \label{eq:hopf-coefficient}
\end{equation}
With the baseline parameters, \(2b-1-\varepsilon b^2=0.5488\). The coefficient is positive at both crossings in these coordinates.

\begin{table}[H]
  \centering
  \caption{Hopf classification at \(a=0.7\), \(b=0.8\), and \(\varepsilon=0.08\). The determinant is \(0.075904\) at both crossings.}
  \label{tab:hopf}
  \begin{tabular}{@{}rrrrl@{}}
    \toprule
    \(I_H\) & \(v_H\) & \(dT/dI\) & \(\ell_1\) & Side with small unstable orbit\\
    \midrule
    0.331281 & \(-0.967471\) & \(+1.631486\) & \(+0.262432\) & \(I<I_H\), nearby\\
    1.418719 & \(+0.967471\) & \(-1.631486\) & \(+0.262432\) & \(I>I_H\), nearby\\
    \bottomrule
  \end{tabular}
\end{table}

In the local amplitude equation \(r'=\mu r+\ell_1r^3+\cdots\), \(\mu\) has the sign of the trace near the crossing. For \(\ell_1>0\), a small nonzero radius exists for \(\mu<0\) and repels nearby trajectories. Thus both crossings are \emph{subcritical Hopf bifurcations}. The calculation determines the stability and location of the small orbit close to each crossing. The large firing loops are investigated separately in the simulations. Subcritical Hopf behaviour in the FitzHugh--Nagumo equations is also discussed by Izhikevich \cite{izhikevich2000}.

\clearpage

\section{Numerical methods}

We solved \eqref{eq:equilibrium} by bisection over \([-3,3]\) for the currents tested, then substituted \((v_*,w_*)\) into both differential equations to check the residuals. We integrated the ODEs with SciPy's \texttt{solve\_ivp} using DOP853, an adaptive Runge--Kutta method \cite{scipy1153}. The relative and absolute tolerances were \(10^{-9}\) and \(10^{-11}\). The initial comparisons report results every \(0.1\) model time units on \([0,150]\). This output grid specifies where values are recorded, while the solver chooses its internal step sizes.

For the displayed firing traces, we counted upward crossings of \(v=0\), estimating each crossing time by linear interpolation between the adjacent output samples. We used the same crossing rule for the sensitivity study. As a numerical check, tightening the sustained-current run to relative and absolute tolerances of \(10^{-11}\) and \(10^{-13}\) changed the sampled states by at most \(4.39\times10^{-8}\).

\section{Simulation results}

\subsection{Small nudge, single spike, and sustained firing}
\label{sec:runs}

At \(I=0\), the equilibrium is approximately \((v_*,w_*)=(-1.1994,-0.6243)\). All three experiments start with recovery \(w_*\) from this equilibrium. Table~\ref{tab:runs} gives the initial voltage offsets and the applied currents.

\begin{table}[H]
  \centering
  \caption{Initial voltage offsets from the zero-current equilibrium, applied currents, and observed spike counts on \([0,150]\).}
  \label{tab:runs}
  \begin{tabular}{lrrr}
    \toprule
    Experiment & Voltage offset & Current \(I\) & Spikes\\
    \midrule
    Small nudge    & \(+0.05\) & \(0.00\) & \(0\)\\
    Larger nudge   & \(+0.90\) & \(0.00\) & \(1\)\\
    Sustained input & \(+0.05\) & \(0.50\) & \(4\)\\
    \bottomrule
  \end{tabular}
\end{table}

\begin{figure}[H]
  \centering
  \includegraphics[width=\linewidth]{figures/voltage_responses.pdf}
  \caption{Voltage responses on a common scale. The dashed horizontal line is the zero-current resting voltage. The right panel also shows the slower recovery state.}
  \label{fig:voltage}
\end{figure}

The small nudge decays quickly. The larger nudge gives a single large excursion and returns to the same resting voltage. These two runs have identical parameters and applied current, so their difference demonstrates excitability driven by the initial state. The sustained input keeps producing spikes because it changes the rates at every time during the run.

\subsection{Locating the firing threshold}
\label{sec:threshold}

The two zero-current runs differ only in their initial voltage, so excitability is a statement about a boundary between them. We locate it by holding the recovery variable at $w_*$, displacing only the voltage, and bisecting on the displacement; the peak voltage over the run separates the two outcomes without any tuning, because a decaying run stays far below $v=1$ while a spike overshoots well past it.

\begin{table}[H]
  \centering
  \caption{The firing threshold at $I=0$, from 45 bisection steps.}
  \label{tab:threshold}
  \small
  \begin{tabular}{@{}lr@{}}
    \toprule
    Quantity & Value\\
    \midrule
    Critical voltage displacement from rest & $0.5555$\\
    Threshold voltage $v_{\mathrm{thresh}}$ & $-0.6439$\\
    Middle branch of the cubic nullcline at $w=w_*$ & $-0.7863$\\
    Peak voltage $0.001$ below the threshold & $-0.3341$\\
    Peak voltage $0.001$ above the threshold & $+1.5887$\\
    \bottomrule
  \end{tabular}
\end{table}

Table~\ref{tab:threshold} reports the threshold measured at $I=0$, with the recovery variable held at its resting value $w_*$ and a peak above $v=1$ taken as the spike criterion. At two test displacements, $0.001$ below and $0.001$ above the estimated threshold, the peak voltage is $-0.3341$ and $+1.5887$ respectively, so peak voltage changes sharply over a narrow band of initial conditions. The located threshold sits $0.1424$ voltage units to the \emph{right} of the middle branch of the cubic nullcline at $w=w_*$, which lies at $-0.7863$. The middle branch is the $\varepsilon\to0$ description of the dividing set; at $\varepsilon=0.08$ the measured threshold and that branch differ by a resolvable amount, which is consistent with the scaling evidence in Section~\ref{sec:scaling}. Bisection locates where the spike criterion changes value; it does not by itself establish that the underlying response is discontinuous, nor identify the dividing set.


\clearpage

\subsection{Phase portraits and equilibrium stability}

In the phase plane, the horizontal coordinate is voltage and the vertical coordinate is recovery. At each point, the direction of motion comes from the two rates in \eqref{eq:voltage} and \eqref{eq:recovery}. The arrows in Figure~\ref{fig:phase} are scaled for readability. They show direction, while time plots show the speed of motion.

\begin{figure}[H]
  \centering
  \includegraphics[width=\linewidth]{figures/phase_portraits.pdf}
  \caption{Phase portraits for zero current and \(I=0.50\). The dot in each panel is the equilibrium at the intersection of the nullclines. The red curve shows the trajectory.}
  \label{fig:phase}
\end{figure}

The zero-current trajectory makes a large loop and returns to the equilibrium. With sustained current, the trajectory approaches a repeating loop around the equilibrium. One trip around that loop corresponds to one firing cycle. Both cases have an equilibrium. Their eigenvalues explain why one attracts a small disturbance and the other repels it.

\begin{table}[H]
  \centering
  \caption{Local stability and numerical checks. The initial state is \(0.01\) voltage units from the corresponding equilibrium. The final column gives the Euclidean distance from that equilibrium at time 150.}
  \label{tab:stability}
  \small
  \begin{tabular}{rrrrlr}
    \toprule
    \(I\) & \(v_*\) & \(w_*\) & Eigenvalues & Stability & Final distance\\
    \midrule
    0.00 & \(-1.1994\) & \(-0.6243\) & \(-0.2513\pm0.2119i\) & Stable & \(2.51\times10^{-10}\)\\
    0.50 & \(-0.8048\) & \(-0.1311\) & \(+0.1441\pm0.1915i\) & Unstable & \(1.70\)\\
    1.60 & \(+1.1043\) & \(+2.2554\) & \(-0.1418\pm0.2719i\) & Stable & \(1.41\times10^{-11}\)\\
    \bottomrule
  \end{tabular}
\end{table}

The numerical perturbation checks agree with the eigenvalues at all three currents. At \(I=1.60\), the equilibrium is locally stable again, at a higher voltage than the zero-current resting state. At \(I=0.50\), the positive real parts predict initial growth of a small disturbance. The large, bounded oscillation that follows is a property of the nonlinear system seen in the simulation.

\clearpage

\subsection{Current sweep and behaviour near the crossings}

We varied \(I\) from 0 to 1.75 in increments of 0.025, using the same initial state as the small-nudge experiment in Table~\ref{tab:runs}. For each run, we recorded the minimum and maximum voltage over \([300,600]\), sampled every \(0.2\) time units. Figure~\ref{fig:sweep} compares these extrema with the equilibrium voltage and the largest eigenvalue real part. The dotted lines mark the two analytically classified subcritical Hopf bifurcations, \(I\approx0.331281\) and \(1.418719\).

\begin{figure}[H]
  \centering
  \includegraphics[width=0.83\linewidth]{figures/current_sweep.pdf}
  \caption{The upper panel shows equilibrium voltage and late-time extrema from a common initial state. The lower panel shows the largest eigenvalue real part.}
  \label{fig:sweep}
\end{figure}

We also simulated currents \(0.004\) below and above each bifurcation. Each run started at the corresponding equilibrium recovery value, with voltage nudged toward the middle of the voltage range by either \(0.01\) or \(0.80\). Table~\ref{tab:refined} shows the resulting voltage span, maximum minus minimum, on \([1700,2000]\).

\begin{table}[H]
  \centering
  \caption{Refined simulations near the stability crossings. A span close to zero indicates a trajectory that has nearly settled over the recorded window.}
  \label{tab:refined}
  \small
  \begin{tabular}{rrrr}
    \toprule
    \(I\) & Largest real part & Span after \(0.01\) nudge & Span after \(0.80\) nudge\\
    \midrule
    0.327281 & \(-0.00326\) & 0.00008 & 3.73783\\
    0.335281 & \(+0.00327\) & 3.76057 & 3.76057\\
    1.414719 & \(+0.00327\) & 3.76057 & 3.76057\\
    1.422719 & \(-0.00326\) & 0.00008 & 3.73783\\
    \bottomrule
  \end{tabular}
\end{table}

At the two tested currents on the stable sides, the small nudge nearly settles and the larger nudge keeps oscillating over the late window. This is numerical evidence consistent with coexistence of a stable equilibrium and a firing state for those settings. The finite windows and the tested starting states limit the strength and range of that conclusion.

\clearpage

\subsection{Sensitivity to the recovery rate}

With \(I=0.50\) and the same starting state, we changed only \(\varepsilon\). We simulated to time 450, recorded output every \(0.1\) time units, and discarded spike crossings before time 150. The mean period is the average time between successive retained crossings.

\begin{center}
  \begin{minipage}[c]{0.55\linewidth}
    \centering
    \includegraphics[width=\linewidth]{figures/recovery_sensitivity.pdf}
    \captionof{figure}{Mean firing period at \(I=0.50\).}
    \label{fig:sensitivity}
  \end{minipage}\hfill
  \begin{minipage}[c]{0.42\linewidth}
    \centering
    \captionof{table}{Recovery sensitivity. Periods are in model time units.}
    \label{tab:sensitivity}
    \small
    \begin{tabular}{@{}rrr@{}}
      \toprule
      \(\varepsilon\) & Late spikes & Mean period\\
      \midrule
      0.06 & 7 & 49.157\\
      0.08 & 8 & 39.475\\
      0.10 & 9 & 33.525\\
      \bottomrule
    \end{tabular}
  \end{minipage}
\end{center}

Faster recovery reduced the average interval between spikes across the tested range. The mean period decreased by about \(32\%\) between the first and last settings. This comparison describes one current and three recovery settings.

\subsection{Scaling of the period with the recovery rate}
\label{sec:scaling}

The slow--fast structure of the model predicts how large that effect should be. A firing cycle spends almost all of its time traversing the outer branches of the cubic nullcline, where motion is governed by \eqref{eq:recovery} and therefore proceeds at a rate proportional to \(\varepsilon\); the two transitions between branches are governed by \eqref{eq:voltage} and take \(O(1)\) time. In the singular limit this gives \(T\propto\varepsilon^{-1}\), equivalently \(\varepsilon T\to\text{constant}\). We tested this by widening the sweep by more than an order of magnitude.

\begin{table}[H]
  \centering
  \caption{Mean firing period at \(I=0.50\) over a wider range of recovery rates, measured after discarding the first 1000 time units.}
  \label{tab:scaling}
  \small
  \begin{tabular}{@{}rrr@{}}
    \toprule
    \(\varepsilon\) & Mean period \(T\) & \(\varepsilon T\)\\
    \midrule
    0.005 & 430.745 & 2.1537\\
    0.010 & 226.648 & 2.2665\\
    0.020 & 122.113 & 2.4423\\
    0.040 & \ 67.951 & 2.7180\\
    0.060 & \ 49.157 & 2.9494\\
    0.080 & \ 39.474 & 3.1580\\
    0.100 & \ 33.525 & 3.3525\\
    \bottomrule
  \end{tabular}
\end{table}

The product \(\varepsilon T\) rises steadily from \(2.15\) to \(3.35\), so it has not reached a constant over the range tested. Fitting \(T\propto\varepsilon^{-\alpha}\) on log--log axes gives \(\alpha=0.909\) over \(\varepsilon\in[0.005,0.02]\) and \(\alpha=0.750\) over \(\varepsilon\in[0.06,0.10]\). The exponent is larger over the smaller-\(\varepsilon\) window and closer to \(1\) there, and \(\varepsilon T\) varies by only a factor of \(1.6\) while \(T\) itself varies by a factor of \(13\). Both observations are consistent with the slow--fast prediction \(T\sim C/\varepsilon\) holding asymptotically. They do not demonstrate convergence: the sweep stops at \(\varepsilon=0.005\), and over this range the \(O(1)\) transition time remains a visible fraction of the cycle. At \(\varepsilon=0.08\) in particular the asymptotic regime has not been reached, so the \(32\%\) reduction reported above is a correct measurement taken in a pre-asymptotic regime, and the inverse law should not be quoted as exact at these parameters.

The two tables give slightly different periods at \(\varepsilon=0.08\): \(39.475\) in Table~\ref{tab:sensitivity} against \(39.474\) in Table~\ref{tab:scaling}. The runs use different measurement windows and output sampling --- Table~\ref{tab:sensitivity} integrates to time \(450\), samples every \(0.1\) time units and discards crossings before time \(150\), while Table~\ref{tab:scaling} discards the first \(1000\) time units --- so the two averages are taken over different sets of cycles and locate each crossing on a different grid. The difference of \(0.001\) time units, around \(3\times10^{-5}\) of the period, is the size expected from that alone.


\section{Bonus study of parameter estimation with a PINN}
\label{sec:bonus}

\subsection{What this section is for, and how to read its result}
\label{sec:bonus-frame}

The core report solves the \emph{forward} problem: given parameters, predict voltage. This section
reverses it. Given only noisy voltage readings, and treating the recovery variable \(w\) as
unobservable --- which it is in any real recording --- can the model's parameter be recovered?

We answer it twice, with a physics-informed neural network and with a classical least-squares fit,
and we state the expected outcome in advance. Both methods are given the same governing equations
and the same initial state, including the unmeasurable \(w_0\). The least-squares baseline has only
a single scalar to search, with the true model integrated at every trial value, so on this
synthetic problem it is the stronger of the two by construction. It should win, and it does.

The comparison is therefore not a contest but a \textbf{benchmark}: least squares is the empirical
reference point for this experiment, and the quantity of interest is how close a network that must
simultaneously learn two state trajectories comes to it. A PINN landing within a small factor of
that reference is evidence that the implementation is correct. Section~\ref{sec:bonus-value}
returns to the assumptions the baseline relies on and to the settings in which they would not hold.

\subsection{Inverse problem and synthetic data}

The core simulations use known parameters to predict voltage. Here we reverse the problem and estimate \(a\) from voltage readings alone. We fix \(b=0.80\) and \(\varepsilon=0.08\), supply the current and initial state, and treat recovery \(w\) as unobserved during fitting. The data come from the FHN equations with \(a_{\mathrm{true}}=0.70\), so the parameter error can be measured directly.

We generated clean reference trajectories on \([0,40]\) at intervals of \(0.10\), using DOP853 with relative and absolute tolerances of \(10^{-11}\) and \(10^{-13}\). For the clean training set, we kept every fourth voltage sample, giving 101 observations. For the sparse, noisy set, we chose one positive-time sample per block of ten and added Gaussian noise with standard deviation \(0.03\), giving 40 observations. The two sets sample the same training trajectory.

\begin{figure}[H]
  \centering
  \includegraphics[width=0.70\linewidth]{figures/pinn_training_observations.pdf}
  \caption{The two training datasets. Dots are the supplied observations and the gray curve is the clean numerical reference.}
  \label{fig:pinn-data}
\end{figure}

Training used \(I=0.50\) and \((v_0,w_0)=(-1.15,-0.62)\). Validation used \(I=0.65\) and \((-1.00,-0.55)\). The final test used \(I=0.90\) and \((-1.20,-0.70)\). Only training observations entered either fit.

\clearpage

\subsection{PINN and least-squares fitting}

A physics-informed neural network (PINN) fits observed data while penalizing disagreement with the governing equations \cite{raissi2019}. Our network takes time as input and predicts voltage \(\hat v\) and recovery \(\hat w\). It uses normalized time, three pairs of sine and cosine features, and three hidden layers of 32 units with \(\tanh\) activations. The output takes the form \((\hat v,\hat w)=(v_0,w_0)+(1-e^{-t})N_\theta(t)\), which gives the known initial state exactly at \(t=0\).

At 241 fixed physics times, automatic differentiation supplies the predicted slopes. The residuals are the predicted slopes minus the rates from \eqref{eq:voltage} and \eqref{eq:recovery}, evaluated at the predicted state and fitted \(\hat a\). The objective is
\begin{equation}
  \mathcal{L}
  =\operatorname{mean}_{\rm obs}\!\bigl[(\hat v-v_{\rm obs})^2\bigr]
    +\operatorname{mean}_{\rm grid}\!\bigl[r_v^2+4r_w^2\bigr].
  \label{eq:pinn-loss}
\end{equation}
On a Colab CPU in double precision, we used 400 voltage-only steps, 1800 joint Adam steps, and up to 300 L-BFGS iterations. The data seed was 2026 and the network seed was 7.

The least-squares baseline integrates the FHN model for each trial \(a\) and minimizes squared voltage errors at the same observation times. Both fits start at \(a=0.50\) with the same allowed interval \([0.20,1.20]\).

\subsection{Held-out results from the original benchmark}

For validation and testing, we used each fitted \(a\) in the same numerical ODE solver at the new current and initial state. This checks parameter transfer to another trajectory. The voltage error is the root mean square error (RMSE) against the corresponding clean reference. Relative parameter error is \(100|\hat a-a_{\mathrm{true}}|/|a_{\mathrm{true}}|\).

\begin{table}[H]
  \centering
  \caption{Parameter recovery and final-test errors at \(I=0.90\). Entries shown as zero are rounded to the displayed precision.}
  \label{tab:pinn-test}
  \small
  \setlength{\tabcolsep}{5pt}
  \begin{tabular}{llrrrr}
    \toprule
    Training data & Method & \(\hat a\) & Rel.\ error & Voltage RMSE & Recovery RMSE\\
    \midrule
    Clean & PINN          & 0.699975 & \(0.0036\%\) & 0.000096 & 0.000022\\
    Clean & Least squares & 0.700000 & \(0.00\%\)   & 0.000000 & 0.000000\\
    Sparse, noisy & PINN          & 0.704093 & \(0.58\%\) & 0.015570 & 0.003554\\
    Sparse, noisy & Least squares & 0.701926 & \(0.28\%\) & 0.007343 & 0.001677\\
    \bottomrule
  \end{tabular}
\end{table}

\begin{figure}[H]
  \centering
  \includegraphics[width=0.91\linewidth]{figures/pinn_heldout_comparison.pdf}
  \caption{Final-test voltages and absolute errors for the sparse, noisy fits. The curves are regenerated from the fitted parameters rounded to six decimal places.}
  \label{fig:pinn-test}
\end{figure}

Both methods matched the single reference spike in the final test. On the clean validation trajectory, the fits from noisy training readings gave voltage RMSEs of \(0.011649\) for the PINN and \(0.005491\) for least squares.


\clearpage

\subsection{Controlled comparison of sampling and noise}
\label{sec:controlled}

The original benchmark changed the number of readings and the amount of noise together. We now vary these two factors separately on the same training trajectory. The dense grid uses 100 positive-time samples from \(0.4\) through \(40\), spaced by \(0.4\) model time units. The sparse grid keeps the second and fifth reading from each block of five dense readings, for a total of 40 samples. The known initial state at \(t=0\) is shared by all fits.

For each noise seed \(2026,\dots,2030\), we draw independent Gaussian voltage errors with mean zero and standard deviation \(0.03\) on the dense grid. The sparse dataset uses the same perturbed readings at its selected times. Each density also has a clean control at the same times. This yields two clean datasets and ten noisy datasets. Related work has studied the sensitivity of FHN parameter estimation to observation noise with other neural architectures \cite{rudi2022}.

Both methods use the same readings, bounds, and known initial state. The PINN architecture, loss, collocation grid, optimizer budget, and initialization are held fixed. Each fitted \(\hat a\) is scored by relative parameter error and by voltage RMSE on the final trajectory at \(I=0.90\). That trajectory is used only for scoring.

\paragraph{How the comparisons isolate each factor.}
For noise, subtract the clean error from the noisy error at the same sample count. For sampling, subtract the dense error from the sparse error using the same noise seed. Positive differences indicate an increase in error. We report medians and full ranges across the five noise realizations. The ranges describe the tested realizations.

\begin{table}[H]
  \centering
  \caption{Controlled experiment results on the final test at \(I=0.90\). Noisy conditions show the median and full range in brackets. The two clean datasets contribute one fit per method.}
  \label{tab:controlled-both}
  \small
  \input{pinn_controlled_results}
\end{table}

Both methods recovered \(a\) closely from clean readings. For the noisy datasets, least squares gave smaller parameter errors and lower held-out voltage RMSEs in all ten matched comparisons. At 100 readings, the median RMSE was \(0.002744\) for least squares and \(0.018354\) for the PINN. At 40 readings, the corresponding medians were \(0.002029\) and \(0.022869\). The least-squares fit solves the exact data-generating equations at each trial value of \(a\). The PINN estimates \(a\) while also learning two state curves and balancing its voltage and equation losses. This benchmark compares those two procedures under a known model and fixed training settings.

\clearpage

\subsection{Separate effects of noise and sample count}

\begin{figure}[H]
  \centering
  \includegraphics[width=\linewidth]{figures/pinn_controlled_comparison.pdf}
  \caption{Parameter error and held-out voltage RMSE by observation condition. Each marker is one fit. The horizontal bars show group medians. All held-out trajectories use the same ODE solver with the fitted parameter.}
  \label{fig:controlled-both}
\end{figure}

\begin{table}[H]
  \centering
  \caption{Paired changes in held-out voltage RMSE. Each noisy contrast summarizes five differences as median [minimum, maximum]. Positive values mean error increased.}
  \label{tab:controlled-paired}
  \small
  \begin{tabular}{@{}lll@{}}
    \toprule
    Change & Least squares & PINN\\
    \midrule
    Add noise at 100 readings & \(+0.002744\) [0.001384, 0.006811] & \(+0.018283\) [0.007749, 0.022230]\\
    Add noise at 40 readings & \(+0.002029\) [0.000996, 0.007552] & \(+0.022628\) [0.004908, 0.030079]\\
    Reduce noisy readings to 40 & \(+0.000645\) [$-0.005139$, 0.004808] & \(+0.008018\) [$-0.008281$, 0.012244]\\
    \bottomrule
  \end{tabular}
\end{table}

At a fixed set of observation times, adding noise raised the test error in every trial for both methods. With noisy readings, reducing the sample count increased the error in three of five least-squares pairs and four of five PINN pairs. The effect of the sample count therefore depended on which readings were retained and on their noise values. For example, least squares had a lower group median with 40 noisy readings even though its median \emph{paired change} was positive. The paired change is calculated by subtracting errors seed by seed before taking the median.

These comparisons summarize five realizations of one noise level. They show performance for the specified FHN model, observation times, and fixed initialization. Other noise levels, observation patterns, or initializations would require separate trials.

\clearpage

\subsection{What this section contributes, and where it stops}
\label{sec:bonus-value}

\begin{table}[H]
  \centering
  \caption{What the bonus study delivers, and where each item is evidenced.}
  \label{tab:bonus-map}
  \small
  \begin{tabular}{@{}p{0.27\linewidth}p{0.66\linewidth}@{}}
    \toprule
    Element & Where, and what was actually done\\
    \midrule
    Advanced numerics &
      Automatic differentiation supplies exact derivatives of the network at 241 collocation
      times; a three-stage schedule of 400 data-only steps, 1800 joint Adam steps and up to 300
      L-BFGS iterations with strong-Wolfe line search. The hard-constraint ansatz
      \((v_0,w_0)+(1-e^{-t})N_\theta(t)\) satisfies the initial condition exactly rather than by
      penalty, removing a loss term and its weight.\\[0.35em]
    Machine learning &
      A physics-informed network estimates the unknown parameter as a trainable weight while
      learning the state curves, from voltage observations alone with \(w\) hidden.\\[0.35em]
    Validation &
      Each fitted parameter is scored on \emph{held-out trajectories at currents never seen in
      training}: validation at \(I=0.65\) and a final test at \(I=0.90\), both from different
      initial states, re-integrated with the ODE solver. Both methods reproduce the correct spike
      count on the held-out run.\\[0.35em]
    Controlled experiment &
      The first benchmark confounded noise with sample count. Section~\ref{sec:controlled}
      replaces it with a paired design: two clean datasets and ten noisy ones (two densities
      \(\times\) five noise seeds), sparse readings copied from the dense ones at shared times,
      one fixed network initialisation. Each of the twelve datasets is fitted by both methods,
      giving twenty-four fits, reported in full.\\[0.35em]
    Interpretation &
      The negative result is explained rather than reported: the baseline is identified as the
      empirical reference for this synthetic experiment, so the PINN's distance from it is
      the quantity of interest.\\
    \bottomrule
  \end{tabular}
\end{table}

\paragraph{Proposed follow-up experiments.}
The baseline's advantage here rests on three assumptions: the model is correct, the initial state
is known, and the unknowns are finitely many scalars. Each can be relaxed, and doing so is the
natural continuation of this study. We have not run these experiments; we set them out as the
tests that would determine whether the ordering above survives.

\begin{itemize}[leftmargin=1.4em]
  \item \textbf{Unknown forcing.} If \(I(t)\) is an unknown \emph{function} rather than a
        constant, a second network head could estimate it pointwise, whereas the baseline would
        first have to assume a parametric form. Whether that changes the ordering is untested.
  \item \textbf{Unknown initial state.} Here both methods were given \(w_0\). In a recording it
        is unobservable, and the baseline would have to search over it too. The experiment is to
        withhold \(w_0\) from both and re-measure the gap.
  \item \textbf{Model misspecification.} With a deliberately wrong right-hand side, the question
        is whether the network's physics residual grows in a way that diagnoses the error while
        the baseline absorbs it into the fitted parameter.
  \item \textbf{Expensive forward solves.} For a spatially extended version of these equations,
        the cost of re-integrating a PDE inside the optimiser at every trial parameter would need
        to be compared against the cost of the collocation formulation.
\end{itemize}

\paragraph{Limitations, stated plainly.}
All observations are synthetic. That is deliberate --- the true
parameter is known only for synthetic data and therefore to measure error rather than merely report a fit --- but it means the
study demonstrates method behaviour, not biological fidelity. It uses one noise level, five noise
realisations, one network initialisation and one unknown parameter. Fitting a published
patch-clamp recording, estimating several parameters jointly, and sweeping the noise level to look
for a crossover are the natural continuations, and none of them is addressed here.


\clearpage

\section{Conclusion and limitations}

The FitzHugh--Nagumo model reproduced three distinct responses by changing the initial voltage or the sustained current. Equilibrium and Jacobian calculations explained when small disturbances decayed or grew, and a Lyapunov function built from \eqref{eq:lyapP} attached a number to the decay of the small disturbance, giving a numerically checked region that reaches voltage displacements of \(0.43433\) against a measured threshold of \(0.5555\). The first Lyapunov coefficient classified both stability crossings as subcritical Hopf bifurcations. The refined runs showed that the initial state can matter near either crossing. At the tested current, increasing the recovery rate shortened the firing period; widening that sweep gave a fitted exponent in \(T\propto\varepsilon^{-\alpha}\) of \(0.750\) over the tested range and \(0.909\) over the smallest decade, consistent with the slow--fast prediction \(T\sim C/\varepsilon\) holding asymptotically rather than at these parameters.

The variables are scaled, so comparing the curves to a particular neuron would require calibration against measured voltage, time, and current. The late-time simulations sample selected currents and initial states. A full description of periodic branches and their stability would require additional analysis. The core Jupyter notebook regenerates the simulations and exports the core figures.

In the controlled inverse study, noise increased the held-out error for both methods at both sample counts. The effect of reducing the sample count varied across paired noise draws. Least squares had lower error than the PINN in all ten noisy comparisons under the known FHN equations. The study uses one noise level, five noise draws, and a fixed network initialization. Generalization across initializations, additional unknown parameters, or biological recordings requires separate experiments.

\subsection*{Reproducibility}

The core notebook derives and evaluates \eqref{eq:hopf-coefficient}, regenerates the current-sweep figure, and reproduces the threshold of Table~\ref{tab:threshold}, the numerically checked region of Section~\ref{sec:lyapunov} and the scaling table of Table~\ref{tab:scaling} in its sections 11--13. The executed bonus notebook (\texttt{fhn\_pinn\_ml.ipynb}) contains both benchmarks. The controlled Colab CPU run used PyTorch 2.14.1, SymPy 1.14.0, SciPy 1.16.3, and Matplotlib 3.10.0. Step~12 uses observation seeds 2026--2030 and network seed 7. Its export contains all 24 rows in \texttt{pinn\_controlled\_trials.csv}, the summary table, and the figure in PDF and EPS.

\begingroup
\footnotesize
\setlength{\parskip}{0pt}
\setlength{\bibsep}{0.1em}
\bibliographystyle{unsrt}
\bibliography{references}
\endgroup

\end{document}
